Assuming that problems can be modeled entirely using polynomials allows the application of a generically convergent technique \cite{Nie2014} that provides easily verifiable certificates of optimality. In this context, we will focus on models where the objective function is a polynomial and the compact feasible set is a finite intersection of polynomial superlevel sets, known as a basic semialgebraic set.

A basic semialgebraic set \( S(g_1, \dots, g_k) = \{\bx \in \R^n \mid g_i(\bx) \ge 0 \; \forall i=1,\dots,k\} \) is defined by polynomials \( g_i \) for some positive integers \( k \) and \( n \). We further assume that the set has the Archimedean property, meaning it is certifiably compact (e.g., one of the \( g_i(\bx) \) takes the form \( a - \|\bx\|^2 \) for some positive \( a \)).

Let $f$ be a polynomial and $Z$ be a basic semialgebraic set. We are interested in finding the global minimum of the following optimization problem:
\[
    \begin{aligned}
        \min_{\bx \in Z} \quad & f(\bx).
    \end{aligned}
\]

Two complementary techniques form the basis of polynomial optimization methods called the Moment-SOS hierarchies \cite{Lasserre2015}. The first approach, which will later form the inner-approximation, is to maximize a lower-bound $\alpha$ on the condition that $f-\alpha$ is non-negative on a semialgebraic set $Z$:
\[
    \begin{aligned}
        \max_{\alpha \in \R}\quad&\alpha\\
        \text{\normalfont s.t.}\quad&f(x) - \alpha \ge 0 \quad \forall x \in Z.
    \end{aligned}
\]
The future outer-approximations will be found while attempting to minimize the Lebesgue integral of $f$:
\[
    \begin{aligned}
        \min_{\mu\in \calP(Z)}\quad&\int_Z f \dif \mu.
    \end{aligned}
\]
where $\calP(Z)$ is the set of regular Borel probability measures on $Z$. The two presented programs are exact convex reformulations of the original. Their practical versions are semidefinite programs which are their strenghtenings and relaxations respectively.

\subsection{Existence of Representing Measures}
The integral of a polynomial with respect to a measure can be written as a function on the moments of the measure. With a multi-index
$\balpha=(\alpha_1,\dots,\alpha_{m})\in\N^{m}$ we use the standard notation $\bx^{\balpha}=x_1^{\alpha_1}\dotsb x_{m}^{\alpha_{m}}$ to define monomials. The $\balpha$-th \emph{moment} of a measure $\mu\in\calP(S)$ is the real number
\[
\mu'(\balpha)=\int_{S} \bx^{\balpha} \dif \mu.
\]
The infinite sequence $\bmu'=(\mu'(\balpha))_{\balpha\in\N^{m}}$ is the \emph{moment sequence} and $\mu$ is called a \emph{representing measure}. The moment problem is the characterization of those real sequences $\by=(y_{\balpha})_{\balpha\in\N^{m}}$ which are moment sequences, that is, $\by=\bmu'$ for some $\mu$. This problem can be decided by checking the positivity of moment expressions via \emph{moment} and \emph{localizing} matrices.

For any $d\in \N$, we define the set of multi-indices truncated to degree $d$ as
\[
    \N^{m}_d=\{\balpha \in \N^{m} \mid |\balpha|\le d\},
\] 
where $\abs{\balpha}=\alpha_1+\dots + \alpha_{m}$. Let $\by=(y_{\balpha})_{\balpha\in\N^{m}}$ be a real sequence and $d\in \N$, the \emph{moment matrix of order $d$} is the matrix $M_d(\by)$ with rows indexed by $\balpha\in\N^{m}_d$ and columns indexed by $\bbeta\in\N^{m}_d$ such that
\[
(M_d(\by))_{\balpha\bbeta}=y_{\balpha+\bbeta}.
\]
The polynomials $g_{i}$ defining the semialgebraic set $S(g_1,\dots,g_k)$ can be written as 
\[
g_{i}(\bx)=\sum_{\bgamma}c_{\bgamma}\bx^{\bgamma}
\] 
for some real coefficients $c_{\bgamma}$, where only finitely many of them are nonzero. Then, the~\emph{localizing matrix of order} $d$ is the matrix $M_d(g_{i},\by)$ with rows indexed by $\balpha\in\N^{m}_d$ and columns indexed by $\bbeta\in\N^{m}_d$, where
\[
(M_d(g_{i},\by))_{\balpha\bbeta}=\sum_{\bgamma}c_{\bgamma}y_{\balpha+\bbeta+\bgamma}.
\]

When a real symmetric matrix $A$ is positive semidefinite, we write $A\succeq 0$. Moment sequences are characterized by positive semidefinitness of moment and localizing matrices \cite[Theorem 2.44(b)]{Lasserre2015}. Specifically, a real sequence $\by=(y_{\balpha})_{\balpha\in\N^{m}}$ with $y_\mathbf{0}=1$ is a moment sequence with a Borel representing measure supported in $S(g_i,\dots,g_k)$ if, and only if,
\begin{equation}\label{moment}
    \begin{aligned}
M_d(\by) &\succeq 0 \quad \forall d \in \N\\
M_d(g_{i}, \by) &\succeq 0 \quad \forall d \in \N,\,\forall k\in K.
    \end{aligned}
\end{equation}

\subsection{Nonnegative Polynomials Using Sums of Squares}
Our next goal is to establish a necessary condition for the nonnegativity of polynomials over a basic semialgebraic set that can be verified algorithmically. Observe that if a polynomial can be decomposed into a sum of squared terms (SOS), it is nonnegative. Further notice, that for any point \(\bx\) in \(S(g_1, \dots, g_k)\), the expression \(g_i(\bx) s(\bx)^2\) is also nonnegative for any polynomial \(s\) and for all \(i = 1, \dots, k\). Quadratic modules extend this concept by summing such expressions. Specifically, the \emph{quadratic module} \(Q(G)\) generated by the set \(G = \{g_1, \dots, g_k\}\) is defined as
\begin{equation}\label{def:QM}
    Q(G) = \left\{\, s_0^2 + \sum_{i\in K} g_is_i^2 \;\middle|\; \text{$s_0,s_i$ are polynomials} \; \forall i\in K=1,\dots,k \,\right\}.
\end{equation}

\begin{theorem}[Putinar's Positivstellensatz \cite{Putinar93}]\label{thm:Putinar}
    If a polynomial $p(\bx)$ is strictly positive for all $\bx \in \R^n$ in $S(g_1,\dots,g_k)$, then $p \in Q\left(\{g_1,\dots,g_k\}\right)$.
\end{theorem}
To formulate the sum of squares condition using semidefinite constraints, we define the basis of the linear space of polynomials with degree at most $d$ as
\[
    [\bx]_d=(\bx^{\balpha})_{\balpha\in \N^{m}_d}.
\]
Any polynomial $p(\bx)$ with the coefficient vector $\mathbf{c}=(c_{\balpha})_{\balpha\in \N^{m}_d}$ can then be expressed as 
\[
    p(\bx) = \sum_{|\balpha|\le d} c_{\balpha} \bx^{\balpha} = \mathbf{c}^\intercal[\bx]_{d}.
\]

A polynomial $p(\bx)$ of degree at most $2d$ is SOS if, and only if, there exists a positive semidefinite matrix $G$ such that $p(\bx)=[\bx]_d^\intercal G [\bx]_d$ \cite[Proposition 2.1]{Lasserre2015}. We can therefore formulate the problem of checking whether such $p(\bx)$ is SOS as the semidefinite feasibility problem for a matrix $G$ indexed by $\bbeta$ and $\bgamma$:
\begin{equation*}
    \begin{aligned}
       \sum\limits_{\bbeta+\bgamma=\balpha} G_{\bbeta\bgamma} &= c_{\balpha}  \quad \forall \balpha, |\balpha|\le 2d,\\
G &\succeq 0.
    \end{aligned}
\end{equation*}
Membership of $p(\bx)$ in $Q(S)$ can then be formulated as a single semidefinite program by applying such conditions to each polynomial variable in Theorem \ref{thm:Putinar} \cite[Chapter 2.4.2]{Lasserre2015}. 


\subsection{Hierarchy of Semidefinite Relaxations}
The final step is to put bounds on the sizes of moment sequences and degrees of polynomials using the Lasserre's hierarchies \cite[Chapter 6]{Lasserre2015}. Each relaxations in the hierarchy is a convex semidefinite program that approximates the solution of the original program. 

We consider the quadratic module $Q_h(G)$ truncated to a degree $h$:
\[
Q_h(G) = \left\{\, s_0 + \sum_{i\in K} g_is_i \;\middle|\; {}
   \begin{aligned}[c]
    & s_0,s_i \; \text{\textsc{sos} polynomials}, \\
    & \deg(g_is_i)\le h, \forall i\in K, \\
    & \deg(s_0)\le h 
   \end{aligned}\,\right\}
\]
where $G = \{g_i \mid i \in K = 1,\dots,k\}$. Instead of testing polynomial positivity by checking its membership in $Q(S)$ the $h$th level of the hierarchy tests for membership in $Q_h(S)$. Similarly, instead of requiring that the constraints hold for all moments, the $h$-th program in a hierarchy only constrains the moment and localizing matrices of order $h$. Moreover, we consider only truncated moment sequences $\hat{\by}$, which are just finite-dimensional real vectors. 

Each level of the hierarchy is then a semidefinite program and by going to higher levels, the original problem can be approximated arbitrarily well.